%% 第 4 章  自动生成于 scripts/tex/md2tex.py，请勿直接手改
\chapter{基于代码索引泄露的触发式随机化}
\label{chap:4}

\section{引言}

第 3 章把持续随机化扩展到了内核 built-in 代码上，并给出了它的适用上界：一次调用只要有一端落在随机化区域内就要经过跳板，随机化不到 1\% 的函数已使常驻开销达 1.98\%，覆盖 40\% 内建代码则达 28.83\%。随机化因此只能限定在一个子集内。然而信息泄露却可能发生在内核的任何位置。本章由此提出只随机化信息泄露的函数：既压低跳板的开销，又把安全防护的覆盖面扩展到全内核。

随机化区域的当前地址集中在跳转桩的分支立即数中，这一处位置同时决定两件事：攻击者能否获知布局，以及系统所付出的常驻开销。本章对它施加只执行内存，使读取行为成为可观测事件；再以该事件决定哪些函数需要被随机化，从而把随机化范围由编译期静态选定改为运行时由泄露事件触发。保护代码索引的机制同时就是识别信息泄露的机制。

本章的组织如下。第 4.2 节从跳板开销、触发策略与保护整个代码区域的代价三条线说明研究背景并归纳挑战；第 4.3 节给出系统设计，包括总体架构、只执行代码区域的选取、信息泄露的识别、以及两级跳板与跳板激活；第 4.4 节给出 x86 与 ARM 两条实现路径与两级跳板的代码生成；第 4.5 节评估；第 4.6 节给出安全性分析；第 4.7 节小结。

\section{研究背景与动机}

\subsection{跳板机制在开销上的局限}

持续随机化方案普遍以某种间接层承载代码索引：Adelie\citep{nikolaev2022adelie} 用只读的全局偏移表，Chaos\citep{chang2024chaos} 与第 3 章用跳板。间接层的代价按经过次数计，而经过次数由随机化范围决定。设随机化函数占全部函数的比例为 p，一次调用只要有一端在随机化区域内就要经过间接层，经过的调用占全部调用的比例为 2p − p$^{2}$，小范围时约为 p 的两倍。每纳入一个函数随机化，被拖上间接层的调用约相当于两个函数的量，因为调用的两端各付一次。第 3.5.5 节的实测给出了这一杠杆在真实内核上的形式：随机化 483 个函数时常驻开销为 1.98\%，扩到 3,285 个函数时为 28.83\%。这构成把随机化推广到整个内核的核心阻碍，也是现有工作局限于部分随机化的原因。

部分随机化的代价是覆盖面。攻击者要利用的漏洞与要复用的代码可以在内核的任何位置，一个在编译期选定的子集只能覆盖它自己。要在不付出全内核跳板开销的前提下把覆盖面扩到全内核，随机化的对象就不能在编译期定死，而要在运行时由泄露事件确定。

\subsection{随机化的触发策略}

固定时间间隔是最常见的触发策略，Shuffler\citep{williams2016shuffler} 采用 50 ms，Adelie 采用 20 ms。它把代价平摊在整个运行期，而收益只在极窄的窗口内产生：在没有攻击的绝大部分时间里随机化是纯粹的开销，在攻击发生的那一刻间隔是否足够短则取决于运气。

事件驱动的策略把触发与程序行为关联。TASR\citep{bigelow2015tasr} 在成对的输入输出系统调用之间触发，依据是任何输出类系统调用都可能被用于泄露、任何输入类系统调用都可能被用于上传载荷；ReRanz\citep{wang2017reranz} 以数据传输的累积量作为触发条件。二者关联的都是可能被用于泄露的行为而非实际发生的泄露，以系统调用或 \lstinline[breaklines=false]@copy_to_user@ 一类的接口作间接代理，对 I/O 密集型负载触发过频，对不经这些接口的泄露则不触发。

探测驱动的策略以检测到的探测行为触发。SafeHidden\citep{wang2019safehidden} 检测对用户态进程中信息隐藏安全区的试探并在触发后重随机化该区域，其保护对象是数据而非代码；MProbe\citep{li2022mprobe} 检测对用户态代码的地址探测与内容探测并收回被探测代码的执行权限；RandFun\citep{li2024randomize} 在检测到探测后随机化被探测函数内部的代码块。这一路线与本章的触发哲学一致，但三者的保护对象都在用户态，且都回避了搬动函数地址：或者保护的是数据，或者只收回权限，或者只重排函数内部而把函数留在原处。

上述策略共同的不足在于触发判据与信息泄露之间只有间接关联。本章的观察是，持续随机化的基础设施本身就可以直接识别信息泄露：随机化区域的当前地址只存在于跳转桩这一处，对它的数据读取就是泄露本身，不需要任何代理。据此可以精确确定需要随机化的函数，即那些跳转桩被读过的函数。

\subsection{保护整个代码区域的代价}

只执行内存用于防御信息泄露时，现有方案默认代码区不被读取。这在用户态成立，进程不读自己的代码；在内核中不成立，内核持续地读取自己的代码：动态跟踪设施在插入与恢复探测点时读取原指令，指令替换与静态键在启动期读取并改写代码，模块加载器与实时补丁在应用重定位时读取代码，性能分析工具在栈回溯与符号解析时读取指令流。防护式的内核只执行内存因此必须为这些合法读取者维护一份豁免清单，漏掉任何一条都会使内核在该路径首次执行时崩溃。

两项内核只执行内存工作给出了两种让步方式。kR\textasciicircum{}X\citep{pomonis2019krx} 的处置是永久性的局部豁免：它克隆 \lstinline[breaklines=false]@get_next@ 与 \lstinline[breaklines=false]@peek_next@ 家族的七个例程以及 \lstinline[breaklines=false]@memcpy@、\lstinline[breaklines=false]@memcmp@、\lstinline[breaklines=false]@bitmap_copy@ 共十个函数，使克隆版本不受其代码读取约束的插桩保护，再以约三百三十行补丁把跟踪与调试设施改接到克隆版本上；其论文在讨论向新内核版本移植时进一步说明，eBPF 即时编译器与实时补丁两项设施需要单独处理。ARM 上以硬件调试观察点实现内核只执行内存的工作\citep{jang2019revisiting}的处置是时间性的整体豁免：内核中存在读取节内字面量池以查询符号信息的合法路径，由于字面量池位于代码节内，对它的读取与攻击者的代码读取在硬件看来完全相同，该工作在这些路径执行期间翻转观察点控制寄存器的使能位，把保护整体关闭。

两种让步的形式不同，根源相同：只要保护对象是整个代码区，就必然要为内核自身的读取者开口子，而这份清单既随内核版本变化，又必须完备。安全性与正确性被绑在了一起。

\subsection{研究动机与挑战}

上述三条线各自指向一个挑战。其一，如何使只执行保护不依赖一份必须完备的豁免清单，这决定机制能否在真实内核上部署。其二，如何把读取事件转化为可靠的随机化触发判据，并使从触发到完成的延迟短于攻击者可利用该次读取的时间，这决定触发式随机化是否真的比固定间隔更及时。其三，如何使常驻开销与随机化范围的规模解耦，这决定覆盖面能否扩大到全内核而开销不随之上升。

\section{系统设计}

\subsection{总体架构}

图~\ref{fig:4-1} 给出本章机制的总体结构：跳转桩区受只执行保护，读取事件驱动跳板激活与随机化；候选集合可扩至全内核，激活集合只含被读过的函数。它建立在第 3 章的三区域布局之上，改动集中在两处：跳转桩区被置于只执行内存保护之下，对它的数据读取产生硬件事件；跳板改为两级，每个候选函数在未被读取时停留在未激活态，只付两条直接分支的代价，被读取后激活为完整的跳板并纳入随机化。

\begin{figure}[!htbp]
    \centering
    \includegraphics[width=\textwidth,height=0.46\textheight,keepaspectratio]{schematic/ch4-architecture-mono}
    \bicaption{\enspace 本章机制的总体架构}{\enspace Overall Architecture of the Proposed Mechanism}
    \label{fig:4-1}
\end{figure}

机制的运行分三步。候选集合在编译期确定，可以覆盖全部可搬移的 built-in 函数；候选函数的一级跳板在未激活态只是一条到跳转桩的直接分支。对跳转桩的任何数据读取被硬件捕获，事件处理记录被读的函数、通知随机化线程、放行本次读取。随机化线程把被读函数激活：改写其一级跳板使之转向二级跳板，二级跳板承担第 3 章跳板的全部职责，该函数随后被随机化，攻击者读到的值失效。被激活的函数在一段时间内未再被读取后降回未激活态。

\subsection{只执行代码区域的设计}

保护对象取跳转桩区而非整个代码区。经第 3 章的处理，随机化区域的当前地址在内存中只存在于跳转桩的分支立即数中：函数内引用是相对的，跨函数引用指向位置固定的跳板，函数指针的值就是跳板地址，栈上的返回地址是与随机化无关的标记。读取跳转桩区等于一次性获得整个区域的布局，除此之外没有别的途径。

正常执行到达函数体的路径全部由直接分支构成：一级跳板是一条到跳转桩的直接分支，跳转桩是一条到函数体的直接分支，两者都只被取指，不含任何数据访问；返回标记的解析同样只经重入桩的 PC 相对计算完成。内核的合法读取者读的都是函数体的指令字节，而跳转桩区里没有任何函数体，也没有任何返回地址落在其中，栈回溯与符号解析不会读到它。因此落在跳转桩区上的任何数据读取都没有合法解释，不需要判定读取者的身份，也不需要为任何读取者维护豁免清单。硬件对取指与取数的区分使该判据可由硬件直接给出：只执行内存的两种实现，x86 的管理态保护键与 ARM 的调试观察点，都只作用于数据访问而对取指透明，第 4.4 节分别说明。第 4.2.4 节的挑战一由保护对象的选取本身解决。

覆盖边界随之确定。本章不覆盖对代码区本身的读取：攻击者可以绕开跳转桩，直接扫描随机化区域的代码收集函数内部的片段，这类读取不落在跳转桩区因而不产生事件。对这类读取，本章提供的保护来自随机化本身：所收集片段的地址在该函数下一次被搬移后失效，而失效之后不存在把旧地址翻译回新位置的路径。攻击者经内核自身导出的地址接口或微架构侧信道获得地址的情形不在本章的威胁模型之内。

\subsection{信息泄露识别}

\textbf{事件的处置。} 处置只有三步，记录、通知、放行。记录的内容是被读取的槽位与对应的函数；通知交给随机化线程，由它决定何时执行激活与切换；放行使读取者得到当前值，而该值在切换完成后即失效。三步都在异常处理路径上完成，不涉及页映射改写，也不为被读函数准备任何副本。放行是有意的：攻击者得到的是一个随即失效的值，而放行使机制不必区分读取者身份，也不会因误判而阻断任何执行。

\textbf{不再需要控制流审计。} 攻击者使用过期地址的行为已由旧地址失效处理：随机化完成后，被搬移函数的旧字节被覆盖为软件断点，触发即被判定为攻击并终止。本章不需要再对它做一次判定，也就不需要为被读函数准备正常副本与陷阱页，那一层机制及其自身引入的攻击面一并省去。

\textbf{被读函数集合的生命周期。} 一次落在跳转桩区的数据读取把该桩对应的函数加入集合，这一步的代价是一次表项写入。集合决定下一次切换重排哪些函数：触发时不对整个候选集合重新布局，只重排集合中的函数，单次随机化的代价因而与被读的规模挂钩，与候选集合的整体规模无关。一次随机化之后，集合中的函数已经迁移到新地址，攻击者此前读到的值随之失效；集合中的函数在一段时间内未再被读取后离开集合并降回未激活态。离开机制不可省略：若集合只进不出，它会随时间单调增长并最终逼近整个候选集合，一次触发的代价退化为全区域重排，激活集合与候选集合的分离也随之消失。

\textbf{失效窗口。} 从读取发生到切换完成之间，攻击者手中的值仍然有效，该窗口的长度即第 4.5.3 节测得的触发延迟。窗口之后，旧地址落到被覆盖的软件断点并被终止。窗口长度是本章的主要安全指标。图~\ref{fig:4-2} 给出事件处置与失效窗口的时序：读取被放行，所得的值在切换完成后失效，越过窗口之后再使用该值会落到被覆盖的软件断点。

\begin{figure}[!htbp]
    \centering
    \includegraphics[width=\textwidth,height=0.46\textheight,keepaspectratio]{schematic/ch4-event-window}
    \bicaption{\enspace 跳转桩读取事件的处置与失效窗口}{\enspace Handling of a Jump-Stub Read Event and the Invalidation Window}
    \label{fig:4-2}
\end{figure}

\subsection{两级跳板与跳板激活}

第 3 章的跳板对每一次跨区域调用约 53 个周期，其中约 36 个周期来自活跃计数的进出，而计数只对正在被随机化的函数有意义。本节据此把跳板分为两级，使未被读取的函数不计数。

\textbf{两级跳板的结构。} 所有调用点调用一级跳板。一级跳板有两种状态，指令序列如下。

\noindent\begin{minipage}{\linewidth}
\begin{lstlisting}[,numbers=none,basicstyle=\ttfamily\scriptsize,keepspaces=true]
; 未激活态：一级跳板只是一条到跳转桩的直接分支，不记账
fixed_in_F:
    b       stub_F                    ; 跳转桩再以 b F_entry 到达 F 的当前位置；F 直接返回原调用点

; 激活态：一级跳板改为一条到二级跳板的直接分支
fixed_in_F:
    b       fixed_in2_F               ; 同一条指令，只改立即数

; 二级跳板与第 3 章的入口跳板完全相同：读当前任务的内外标志，标志为 0 则计数加一、
; 置标志并查阻断标志，标志为 1 则只查阻断标志；经同一个跳转桩调用 F；返回时按进入时
; 保存的标志值决定是否减一
fixed_in2_F:
    ...                               ; 第 3 章 fixed_in 的序列，见第 3.3.2 节

; 出口跳板与第 3 章相同：标志为 0 的执行流不记账，直接调用并按标记返回
fixed_out_G:
    ldr     w17, [x_cur, #inside]
    stp     x17, x30, [sp, #-16]!
    cbz     w17, .Lpassive            ; 标志为 0：未经计数进入，不记账也不查白名单
    ...                               ; 标志为 1：第 3 章的出口跳板序列
.Lpassive:
    bl      G
    ldp     x17, x30, [sp], #16
    b       resolve_token             ; 按标记解析到 F 的当前位置
\end{lstlisting}
\end{minipage}

两种状态共用同一对跳转桩。返回标记的解析总是经重入桩完成，与函数处于哪种状态无关，解析逻辑不分叉；被保护的对象仍只有跳转桩一处，未激活函数的当前地址同样只存在于这里，读取它同样产生事件。

\textbf{出口跳板的记账判据。} 第 3 章的出口跳板在离开区域时对活跃计数减一、返回时加一，前提是执行流进入区域时曾加一。未激活函数经一级跳板进入时不加一，若出口跳板不加区分地减一，计数就会失衡。第 3 章的跳板已经以每任务的内外标志记录一条执行流是否被计数：从区域外经入口跳板进入时置 1，离开时清零。未激活函数经一级跳板进入时不碰标志，其后的出口跳板读到标志为 0，本次离开不记账，也不查白名单，直接调用外部函数并按标记返回。本章因此不需要改动第 3 章的任何一条跳板序列，二级跳板就是第 3 章的入口跳板。这一判据在任意嵌套下自洽：激活函数 F 经出口跳板调用外部函数 G 时标志已清零，G 再调用一个未激活函数 H，H 的离开与返回都不记账，G 返回到 F 的出口跳板时标志置 1、计数加一，F 返回到二级跳板时清零、减一。未激活函数 H 调用激活函数 F 时，二级跳板读到标志为 0，按从区域外进入处理并计数，F 因而受到保护。激活前已经进入的执行流标志为 0，其后续的离开同样不记账，因此激活的时机不影响计数的配平。中断、软中断与异常处理程序在入口把标志保存并清零，作为独立的执行流记账（第 3.3.4 节），处理程序经未激活跳板进入区域再向外调用时出口跳板同样读到 0。出口跳板读到 0 只有一个含义，即这条执行流进入区域时未被计数，来源只有未激活的一级跳板。

\textbf{周期分析。} 表~\ref{tab:4-1} 按指令序列给出两级跳板各段的周期数。沿用第 3 章的各段由该章微基准实测的时长按测量机时钟折算；本章新增的两段在关抢占关中断的紧循环中单独计时，口径与第 3 章表~\ref{tab:3-3} 相同，真实规模下的单次调用开销另见第 4.5.5 节。未激活态的入口只有一级跳板与跳转桩两条直接分支；未激活态的出口只有标志读取、条件分支、调用与标记解析。激活态的入口比第 3 章多一条直接分支，出口与第 3 章相同。

\begin{table}[!htbp]
    \bicaption{\enspace 两级跳板各段的指令序列与周期分析}{\enspace Instruction Sequence and Cycle Analysis of the Two-Level Trampolines}
    \label{tab:4-1}
    \centering
    \scriptsize
    \setlength{\tabcolsep}{3pt}
    \renewcommand{\arraystretch}{1.25}
    \hyphenpenalty=10000\exhyphenpenalty=10000\emergencystretch=1em
    \begin{tabular}{>{\raggedright\arraybackslash}p{\dimexpr 0.0750\textwidth-2\tabcolsep\relax}>{\raggedright\arraybackslash}p{\dimexpr 0.3000\textwidth-2\tabcolsep\relax}>{\raggedright\arraybackslash}p{\dimexpr 0.4250\textwidth-2\tabcolsep\relax}>{\raggedleft\arraybackslash}p{\dimexpr 0.1100\textwidth-2\tabcolsep\relax}>{\raggedleft\arraybackslash}p{\dimexpr 0.0700\textwidth-2\tabcolsep\relax}}
        \hline
        状态 & 段 & 指令 & 静态指令数 & 周期\\
        \hline
        未激活 & 入口：两条直接分支 & \lstinline[breaklines=false]@b stub_F@ / \lstinline[breaklines=false]@b F_entry@ & 2 & 约 1\\
        未激活 & 出口：标志判定后直接调用与标记解析 & \lstinline[breaklines=false]@ldr@ / \lstinline[breaklines=false]@stp@ / \lstinline[breaklines=false]@cbz@ / \lstinline[breaklines=false]@bl G@ / \lstinline[breaklines=false]@ldp@ / \lstinline[breaklines=false]@cbz@ / \lstinline[breaklines=false]@b resolve_token@ & 7 & 约 4\\
        未激活 & 一次往返合计 &  & 9 & 约 5\\
        激活 & 入口：一级跳板的直接分支 & \lstinline[breaklines=false]@b fixed_in2_F@ & 1 & 约 1\\
        激活 & 入口：标志读写、计数进出、经跳转桩调用 & 第 3 章入口跳板序列 & 18 & 约 37\\
        激活 & 出口：标志判定、计数进出、白名单、调用与解析 & 第 3 章出口跳板序列 & 约 30 & 约 16\\
        激活 & 一次往返合计 &  & 约 49 & 约 54\\
        \hline
    \end{tabular}
\end{table}

未激活态一次往返约 5 个周期，激活态约 54 个周期，相差一个数量级。常驻开销由此与候选集合的规模解耦：候选集合可以扩到全内核，每次调用只为激活集合中的函数付完整的跳板代价，而激活集合只含被读过的函数。第 4.2.4 节的挑战三由此解决。

\textbf{激活。} 激活只改写一级跳板的那一条直接分支，把它的目标由跳转桩改为二级跳板，降级则反向改回。激活是一次 4 字节的原子写入，代价与该函数的调用点数量无关，第 3 章改写条目数等于函数个数这条不变式在激活操作上同样成立。激活之后，新进入的执行流全部经二级跳板计数，随机化线程按第 3 章的流程搬移该函数并改写跳转桩。

\textbf{激活后旧副本的回收。} 激活之后，该函数在激活前的执行流仍可能停留在旧副本中，旧副本要等它们全部离开才能填陷阱。这批执行流不需要逐个跟踪，返回标记替它们完成了大部分排空：向外调用与阻塞的执行流在返回时按标记解析，跳转桩已指向新变体，它们自动落到新位置；睡眠或阻塞的执行流一定阻塞在被调用的外部函数里，返回时同样落到新变体。旧副本的滞留者因此只剩在函数体内直线执行、一次外部调用都没发出的执行流，其上界即第 3 章实测的等待清空时间，0.04 至 0.18 μs。残余的被抢占执行流以一次性扫描判定：先激活，使不再有任何合法路径新进入旧副本，再扫描全部内核栈、被切换出去的任务其保存现场以及各核的当前程序计数器，检查是否有落在旧副本区间内的返回地址或程序计数器。扫描通过即填陷阱；扫描发现滞留者则推迟重试，兜底为 \lstinline[breaklines=false]@synchronize_rcu_tasks@，内核回收 ftrace 跳板面对的是同一个问题，用的正是这个原语。兜底的宽限期为毫秒量级，置于关键路径之外。图~\ref{fig:4-3} 给出两级跳板的状态与激活流程：左侧是未激活态与激活态的一级跳板及二者共用的跳转桩，右侧是激活后的四步，改写一级跳板、计数进入、搬移、扫描后回收旧副本。

\begin{figure}[!htbp]
    \centering
    \includegraphics[width=\textwidth,height=0.46\textheight,keepaspectratio]{schematic/ch4-two-level-trampoline}
    \bicaption{\enspace 两级跳板与跳板激活}{\enspace Two-Level Trampolines and Trampoline Activation}
    \label{fig:4-3}
\end{figure}

\textbf{降级。} 被激活的函数在一段时间内未再被读取后降回未激活态：把一级跳板改回到跳转桩的直接分支，此后新进入的执行流不再计数。降级不需要排空，因为未激活的函数不会被搬动。降级的时机由离开策略决定，可选随机化后清空、按时间老化、按容量上限淘汰三种，本章实现采用随机化后清空。

\section{系统实现}

\subsection{x86 平台：基于管理态保护键的只执行内存}

内核之上不存在更高的强制层，以内核自己维护的页表实现的只执行内存可被内核态任意写关闭，且同一物理页在内核中通常同时存在于内核映像映射与直接映射之下，后者是前者的别名映射，若只作用于其中一份映射，攻击者经由另一份即可读到内容。这一问题不必依赖虚拟化层解决。IskiOS\citep{gravani2021iskios} 指出保护键记录在页表项中，只要让两份映射使用同一个保护键，别名即被一并覆盖。本文采用更直接的做法，且在两个平台上一致：把跳转桩区所占的物理页从直接映射中撤去，使其只剩内核映像映射一份，只执行保护只需施加在这一份映射上。撤去直接映射复用内核既有的接口，\lstinline[breaklines=false]@memfd_secret@ 与 \lstinline[breaklines=false]@vmalloc@ 释放可执行内存时用的是同一条路径；x86-64 的直接映射以大页建立，撤去其中的页需要先把所在的大页拆成 4 KB 页表项，跳转桩区连续且对齐，拆分只波及它所在的少数几个大页。本节据此以管理态保护键实现 x86 上的只执行内存。

保护键最初用于用户态的进程内隔离，已被用来以极低的域切换成本在同一地址空间内划分互不信任的部分\citep{vahldiek2019erim,hedayati2019hodor}，其管理态版本近年又被用于内核分区化\citep{guo2025bulkhead,huang2022ksplit}。管理态保护键为每个页表项分配一个键，并由一个每处理器核的寄存器保存各键的访问禁止与写禁止位。关键性质是保护键只作用于数据访问，取指不受其约束，因此把跳转桩区的页标上某个键并置位该键的访问禁止位之后，经跳转桩的直接分支照常取指而跳转桩不可被读取，只执行语义直接成立。

每核寄存器这一点对检测式用法尤其重要。事件处理要放行本次读取，需要短暂地开放读权限。在以二级页表实现的方案下，权限作用于物理页而物理页为所有核共享，放行窗口对其它核同时打开，攻击者可在别的核上趁此读取而不被检测。保护键的权限状态保存在每核寄存器中，开放与恢复都只影响本核，不经页表也不需要跨核失效，上述窗口不存在。更新保护键寄存器需要一次写模型专用寄存器的操作，其代价在百周期量级；稳态下寄存器保持访问禁止，只有在一次读取事件发生时才临时开放一次，该代价只在事件上支付。两个平台共同的映射布局、各自的硬件保护以及改写跳转桩的路径见图~\ref{fig:4-4}。

\subsection{ARM 平台：基于硬件调试观察点的实现}

页表权限在 AArch64 的 EL1 上无法表达只执行：去除内核页的读权限会产生取指中止异常，无论该页由操作系统页表还是由二级页表映射。这一结论针对本文实验平台所属的 ARMv8 与 ARMv9.0 架构，较新的权限间接与权限覆盖扩展改变了页表权限的表达能力。

在不提供只执行权限的嵌入式平台上，已有工作以非特权访存指令配合内存保护单元\citep{kwon2019uxom}或以调试监视比较器\citep{shen2020picoxom}达到同样的语义，后者与本节的思路最为接近。硬件调试观察点提供了另一条途径，且与只执行语义天然契合：观察点只对数据访问触发，取指从被监控区域取出指令不产生任何异常。观察点的配置位于调试寄存器中，独立于内核所维护的页表之外。覆盖范围由地址掩码决定而非由观察点数量决定，\lstinline[breaklines=false]@DBGWCR<n>_EL1@ 的 \lstinline[breaklines=false]@MASK@ 域可掩蔽 3 至 31 位地址，单个观察点最大可覆盖 2 GB 区域，跳转桩区集中放置后一个观察点即可覆盖。用观察点实现内核只执行内存已有先例\citep{jang2019revisiting}，本节复用同一硬件机理，差别在于监控对象与事件处置：既有工作监控整个 text 区间并把观察点异常作为终止信号，因而必须为读取字面量池的合法路径整体关闭保护；本节只监控跳转桩区，字面量池位于代码节内而不在跳转桩区，合法读取者的读取不落在监控范围内，保护始终处于开启状态，观察点异常作为触发信号交给随机化线程。

观察点是每个处理器核独立的调试资源，必须在每个核上分别配置，并在处理器上下线与电源状态切换时重新编程。内核的 \lstinline[breaklines=false]@hw_breakpoint@ 子系统统一管理这些寄存器并向 \lstinline[breaklines=false]@perf@ 的硬件断点与 \lstinline[breaklines=false]@ptrace@ 的硬件观察点提供服务，本机制经该子系统注册一个保留项长期独占其中一个，并在上下线回调中重建配置。掩码机制要求被监控区域按其大小对齐且大小为 2 的幂，这一约束由第 3 章的布局分配满足；由于监控对象只是跳转桩区而非整个代码区，对齐带来的地址空间取整代价很小。观察点作用于虚拟地址区间，跳转桩区所在的物理页在直接映射中另有一份别名，观察点只盯映像映射那段虚拟地址，管不到直接映射那段。这正是两个平台都撤去直接映射的原因：撤去之后跳转桩区只剩被监控的那一份映射，任何数据读取都必经它。arm64 只在 \lstinline[breaklines=false]@rodata_full@ 配置下允许撤去直接映射中的页，该配置是默认值，此时线性映射本就按 4 KB 页建立，不存在拆分大页的代价。

\textbf{跳转桩的改写路径。} 撤去直接映射并把映像映射置为只执行之后，随机化线程自己也失去了写跳转桩的路径，改写经内核修改代码的标准路径 \lstinline[breaklines=false]@text_poke@ 完成：在关中断的条件下为目标页建立一份每核私有的临时可写映射，x86-64 经专用的改写地址空间加 fixmap，arm64 经 \lstinline[breaklines=false]@FIX_TEXT_POKE0@ fixmap，写入立即数后立即撤去该映射。这份映射在其存在期间是跳转桩区的第三份别名，但它只在本核可见、关中断期间存在、寿命为一次写入，且改写发生在活跃集合已空、阻断标志置位的切换阶段，攻击者要利用它必须恰好在这个核、这一段微秒级窗口内发起读取。第 6.2 节把这一临时映射列入本文的残余风险。图~\ref{fig:4-4} 给出两个平台共同的映射布局、各自的硬件保护与改写路径：上半是撤去直接映射之后只执行保护所施加的那一份映像映射，x86 以保护键实现、ARM 以观察点实现，权限状态都在每核寄存器中；下半是改写跳转桩时经 \lstinline[breaklines=false]@text_poke@ 建立的每核私有临时可写映射。

\begin{figure}[!htbp]
    \centering
    \includegraphics[width=\textwidth,height=0.46\textheight,keepaspectratio]{schematic/ch4-pkey-aliasing}
    \bicaption{\enspace 只执行保护施加的映射与跳转桩的改写路径}{\enspace The Mapping That Carries Execute-Only Protection and the Jump-Stub Rewrite Path}
    \label{fig:4-4}
\end{figure}

\subsection{两级跳板的代码生成}

编译器为每个候选函数生成一级与二级两个跳板。一级跳板保留原函数名并承担取地址语义，其内容在两种状态下都是一条直接分支，未激活态指向跳转桩，激活态指向二级跳板，在 AArch64 上为 4 字节的 \lstinline[breaklines=false]@b@，在 x86-64 上为 5 字节的 \lstinline[breaklines=false]@jmp rel32@，激活只改写这一条指令的立即数。二级跳板即第 3 章的入口跳板。跳转桩集中放在跳转桩区，该区域按 2 的幂对齐以满足观察点掩码的要求；一级跳板与二级跳板位于位置固定的跳板段。每个候选函数因此比第 3 章多一个一级跳板，体积代价按候选规模线性增长，其实测值见第 4.5 节。

\section{实验评估}

本节回答四个问题：机制在何种负载下会被内核自身的读取触发、触发之后多快使泄露失效、两级跳板是否真的把常驻开销与候选范围解耦、以及触发延迟是否足以在攻击者利用泄露地址之前使其失效。全部攻击行为由内核内受控自测驱动：本文不引入任何向用户态暴露内核读写的攻击载体，由内核内一段受控代码按设定的策略读取跳转桩，以此实例化攻击者读取跳转桩这一被检测行为。

\subsection{实验配置}

主平台 ARM Neoverse N2（观察点路径），副平台 Intel Xeon Gold 6448Y（管理态保护键路径），均 128 GB DDR5，Intel X710 万兆网卡，Samsung PM9A3 NVMe，Linux 6.8，LLVM 14。负载集合固定为下列两组，全章各实验引用时按此命名。

基础负载组：空载（仅系统守护进程）、LMBench 3.0 全套、内核编译（\lstinline[breaklines=false]@make -j64 defconfig@）、fio（\lstinline[breaklines=false]@randread-4k@、\lstinline[breaklines=false]@seqread-1m@，\lstinline[breaklines=false]@--direct=1@）、netperf（\lstinline[breaklines=false]@TCP_STREAM -m 1024@、\lstinline[breaklines=false]@TCP_RR -r 1,1@）、wrk 对 nginx（\lstinline[breaklines=false]@-t8 -c256 -d30s@）、SPEC CPU2006（\lstinline[breaklines=false]@--size=ref --iterations=3@）。

设施开启组：在内核编译负载上分别开启 ftrace function tracer、在十个热点函数上设 kprobes、\lstinline[breaklines=false]@perf record -e cycles@、挂载一个 eBPF 跟踪程序、加载一个 livepatch 补丁、执行 \lstinline[breaklines=false]@kexec -l@（只装载不切换）。

每个配置 10 次运行，报告中位数与四分位距；延迟类指标另报 p99；每次运行前丢弃首轮预热。

\subsection{两条路径的单次检测开销}

一次代码索引读取所触发的检测开销直接决定了机制在高频读取下的可承受程度。本节在两个平台上分别测量从硬件事件产生到读取指令重新执行完毕的完整路径，并按五个阶段分解：异常入口，即进入处理程序并保存现场；定位槽位与记录，即由出错地址算出被读的跳转桩序号并写入被读函数记录；通知，即把该函数交给随机化线程；权限开放与恢复，即为放行本次读取临时开放权限并在读取完成后恢复；返回，即退出处理程序并重新执行被中断的读取指令。每个平台各测 10 次取中位，结果见表~\ref{tab:4-2}。

\begin{table}[!htbp]
    \bicaption{\enspace 两个平台单次检测开销的分解}{\enspace Breakdown of the Single-Detection Cost on the Two Platforms}
    \label{tab:4-2}
    \centering
    \footnotesize
    \setlength{\tabcolsep}{5pt}
    \renewcommand{\arraystretch}{1.15}
    \begin{tabular}{lrrrr}
        \hline
        阶段 & \multicolumn{2}{c}{ARM 观察点} & \multicolumn{2}{c}{x86 保护键}\\
        \cline{2-3}\cline{4-5}
         & 时长 & 占比/\% & 时长 & 占比/\%\\
        \hline
        异常入口 & 0.31 μs & 23.5 & 0.28 μs & 29.2\\
        定位槽位与记录 & 0.24 μs & 18.2 & 0.24 μs & 25.0\\
        通知随机化线程 & 0.18 μs & 13.6 & 0.18 μs & 18.8\\
        权限开放与恢复 & 0.42 μs & 31.8 & 0.15 μs & 15.6\\
        返回并重执行读取指令 & 0.17 μs & 12.9 & 0.11 μs & 11.4\\
        \hline
        \textbf{合计} & \textbf{1.32 μs} & \textbf{100.0} & \textbf{0.96 μs} & \textbf{100.0}\\
        \hline
    \end{tabular}
\end{table}

两条路径的总开销分别为 1.32 μs 与 0.96 μs，同在微秒量级。差异集中在权限开放与恢复一项：观察点路径需要写调试寄存器并做一次上下文同步，占其总开销的 31.8\%；保护键路径只需两次写模型专用寄存器，一次开放一次恢复，占 15.6\%。其余四项在两个平台上接近，因为它们执行的是同一套软件逻辑，与硬件载体无关。这一分解也解释了为什么两条路径都把权限状态放在每处理器核的部件上：放行窗口只对本核打开，既不需要跨核失效，也不给其它核留下可乘之机。

触发频率与检测到随机化完成的延迟在两个平台上没有实质差别。基础负载组在 x86 平台上的触发频率为每秒 0.02 至 0.41 次，与 ARM 平台的 0.02 至 0.44 次处于同一区间，因为触发来自内核自身的代码读取行为，由软件路径而非硬件载体决定；延迟的中位为 11 至 19 μs、99 分位不超过 102 μs，略低于 ARM 平台，与单次检测开销的差值一致。以下各节的数据除另有说明外均测自主平台。

\subsection{触发频率与触发延迟}

本节使用第 4.5.1 节的负载配置，测试不同负载下触发式随机化每秒触发次数以及从事件到随机化完成的延迟分布。

\begin{table}[!htbp]
    \bicaption{\enspace 各类负载下的触发频率与读取来源}{\enspace Trigger Frequency and Read Sources under Various Workloads}
    \label{tab:4-3}
    \centering
    \scriptsize
    \setlength{\tabcolsep}{3pt}
    \renewcommand{\arraystretch}{1.25}
    \resizebox{\ifdim\width>\textwidth\textwidth\else\width\fi}{!}{%
    \begin{tabular}{lrll}
        \hline
        负载配置 & 触发次数/秒 & 读操作类型 & 读取来源\\
        \hline
        空载 & 0.02 & 符号解析 & 周期性内核线程打印带符号的地址\\
        LMBench & 0.11 & 栈回溯、符号解析 & 系统调用路径上的回溯与地址反查\\
        fio 随机小块读 & 0.35 & 栈回溯 & 块层完成路径的回溯\\
        netperf \lstinline[breaklines=false]@TCP_RR@ & 0.28 & 栈回溯 & 网络收发路径的回溯\\
        内核编译 & 0.44 & 重定位代码 & 模块加载器应用重定位时读取代码\\
        wrk & 0.31 & 栈回溯、符号解析 & Web 服务路径的回溯与地址反查\\
        SPEC CPU2006 & 0.02 & 符号解析 & 几乎不进内核，仅周期线程的地址反查\\
        开启 ftrace & 6.80 & 读指令 & 插入与恢复探测点时读取函数入口的原指令\\
        开启 kprobes & 4.20 & 读指令 & 注册探测点时校验指令边界并读取原指令\\
        perf record & 5.40 & 栈回溯、符号解析 & 采样时的回溯与地址反查读取指令流\\
        eBPF 跟踪 & 3.60 & 读指令 & 即时编译器写入代码页前读取\\
        livepatch & 0.90 & 读指令、重定位代码 & 应用补丁时读取被替换函数的入口\\
        kexec 装载 & 1.20 & 校验映像 & 校验时读取代码段\\
        \hline
    \end{tabular}
    }
\end{table}

表中数值测于监控范围覆盖整个随机化区域与跳板区的配置。在该配置下，基础负载组的触发频率在每秒 0.02 至 0.44 次之间，即平均数秒才发生一次随机化；开启动态跟踪设施后上升一个数量级，ftrace 达每秒 6.80 次，kprobes 达每秒 4.20 次，因为这两类设施在插入探测点时读取函数体的指令，而函数体在该配置下处于监控范围之内。读操作类型一列说明了收窄监控范围为何有效：栈回溯与符号解析读的是函数体内的指令与内核的符号表，读指令与重定位代码读的是函数体的字节，校验映像读的是整个代码段，这些读取的对象都不在跳转桩区内。这一频度决定了机制在开启跟踪设施的环境中的占空比，也是离开策略要面对的负载。

\begin{table}[!htbp]
    \bicaption{\enspace 检测到随机化完成的延迟分布}{\enspace Distribution of the Latency from Detection to Re-randomization Completion}
    \label{tab:4-4}
    \centering
    \footnotesize
    \setlength{\tabcolsep}{3pt}
    \renewcommand{\arraystretch}{1.25}
    \begin{tabular}{lrrrr}
        \hline
        负载配置 & 中位 & p95 & p99 & 最大\\
        \hline
        空载 & 12 μs & 31 μs & 58 μs & 74 μs\\
        LMBench & 15 μs & 38 μs & 71 μs & 96 μs\\
        fio 随机小块读 & 19 μs & 52 μs & 96 μs & 138 μs\\
        netperf \lstinline[breaklines=false]@TCP_RR@ & 17 μs & 45 μs & 88 μs & 121 μs\\
        内核编译 & 21 μs & 58 μs & 110 μs & 152 μs\\
        开启 ftrace & 24 μs & 66 μs & 128 μs & 176 μs\\
        开启 kprobes & 22 μs & 61 μs & 119 μs & 164 μs\\
        受控自测（密集探测） & 18 μs & 49 μs & 92 μs & 127 μs\\
        \hline
    \end{tabular}
\end{table}

表~\ref{tab:4-4} 给出检测到随机化完成的延迟，其数值测于与表~\ref{tab:4-3} 相同的配置。基础负载下中位在 12 至 21 μs 之间，99 分位不超过 110 μs；开启跟踪设施后因触发更频繁而略有拉长，最大值 176 μs。这段延迟就是一次泄露的地址从被读到失效的时间，其安全含义在下一节以竞态实验单独评估。

\subsection{泄露地址的有效期与利用竞态}

第 4.5.3 节测得的触发延迟只说明地址多久失效，不直接回答它是否足够短。要回答这一点，需要把它与攻击者利用一个泄露地址所需的时间放在同一条时间轴上。本节不去估计难以定义的“凑齐 gadget 的时间”，而是让受控自测把从读取跳转桩到控制流真正转到该地址的完整路径跑一遍，直接测量这场竞态的结果。

\textbf{实验构造。} 受控自测扮演能力最强的攻击者：预先在一处由它掌握的内核数据结构中放好一个函数指针，并把利用逻辑除目标地址外全部就绪。一次试验分三步。其一，读取某个随机化函数的跳转桩，从其分支立即数得到该函数的当前地址，这一读取落在只执行监控范围内，触发一次随机化。其二，立即把读到的地址写入那个函数指针。其三，经一条内核路径解引用该指针。若解引用发生在随机化完成之前，控制流转到仍然有效的地址，记为攻击成功；若随机化已经完成，该地址已失效，指针目标落到被覆盖的软件断点并被判定为攻击，记为防御成功。

\textbf{自变量是攻击者的暂存程度，即从读取到解引用之间的时延。} 攻击者暂存得越充分，这段时延越短。本节扫描这段时延，报告每一档下的攻击成功率，成功率由高到低转折的位置就是防御生效的边界。这样得到的是一条可复现的成功率曲线，而非一个含义模糊的“凑齐时间”，其定义只依赖两个有明确物理意义的时刻：读取与解引用。

\begin{table}[!htbp]
    \bicaption{\enspace 泄露地址利用竞态的结果}{\enspace Outcome of the Leaked-Address Exploitation Race}
    \label{tab:4-5}
    \centering
    \scriptsize
    \setlength{\tabcolsep}{3pt}
    \renewcommand{\arraystretch}{1.25}
    \begin{tabular}{lrl}
        \hline
        读取至解引用的时延 & 攻击成功率 & 说明\\
        \hline
        1 μs & 99.6\% & 远快于随机化，单地址单次利用几乎必然成功\\
        10 μs & 71.3\% & 接近触发延迟中位，攻击成功取决于本次随机化快慢\\
        20 μs & 38.5\% & 多数攻击尝试已被随机化抢先\\
        50 μs & 8.2\% & 落在延迟分布尾部之外者才成功\\
        92 μs & 1.0\% & 与 99 分位对齐：99\% 的试验中地址已失效\\
        150 μs & 0.1\% & 超过最大观测延迟，几乎全部失败\\
        \hline
    \end{tabular}
\end{table}

\begin{figure}[!htbp]
    \centering
    \includegraphics[width=0.82\textwidth]{plots/ch4-race-success}
    \bicaption{\enspace 泄露地址利用竞态的成功率随读取至解引用时延的变化}{\enspace Success Rate of the Leaked-Address Race versus the Read-to-Dereference Delay}
    \label{fig:4-5}
\end{figure}

图~\ref{fig:4-5} 中的竖线为第 4.5.3 节触发延迟的中位、99 分位与最大值。从结果来看，多步攻击被可靠阻断。真实的内核代码复用攻击不是单地址单次利用，它需要多个不同函数的地址串成链。每读取一个地址都触发一次随机化，而随机化在读取后 12 至 176 μs 内完成，使先读到的地址在后续地址读齐之前就已失效，攻击者无法同时持有一个足够大的有效地址集合。表~\ref{tab:4-5} 中读取至解引用时延超过 92 μs 时成功率降到 1\% 以下，说明任何需要在读取之后再花费近百微秒去准备的攻击都几乎必然失败，而串联多个地址、构造并布置一条完整的劫持链正属于这一类。

但是，单地址的一次性利用是本章不消除的残余风险。若攻击者把一切都预先暂存好，只等填入一个地址就立即在一条热路径上解引用，读取至解引用的时延可以压到微秒级，此时它快于随机化，表~\ref{tab:4-5} 的 1 μs 一档成功率 99.6\% 正是这一情形。本章的触发延迟无法赢下这场单次竞态，它保证的是这个泄露地址此后失效、攻击者不能第二次使用同一次泄露，而不是保证第一次使用也失败。填入的这个地址是否合法，需要在控制流转移时另作检查，这正是第 5 章对跨区域转移施加指针认证的原因：被写入函数指针的地址即便当下仍然有效，也无法通过调用点上下文的认证。本章使泄露地址短命，第 5 章使被篡改的指针无法通过验证，两者共同覆盖这场竞态的两侧。

\subsection{未激活与激活状态的单次调用开销}

本节选择不同候选集合规模 N ∈ \{483, 3,285, 12,000\}，全部置于未激活态，测试一次跨区域调用的边际开销，以及 L1 指令缓存缺失率、iTLB 缺失率与分支预测失败率。另在同一构建上运行 LMBench \lstinline[breaklines=false]@lat_syscall@ 与 \lstinline[breaklines=false]@lat_ctx@ 取端到端印证。

未激活态的一次进入只经两条直接分支，在 483 个跳板的紧循环下单次 0.51 ns。紧循环指只反复执行被测操作、不含其它工作的微基准循环，其结果不含指令缓存与分支预测器的压力，是未激活态边际开销的下界。但是 12,000 个候选函数意味着 12,000 个一级跳板与跳转桩，约二百 KB 常热的跳板代码与 24,000 条直接分支，指令缓存与分支目标缓冲的压力在紧循环无法准确测试。

表~\ref{tab:4-6} 在真实规模下给出这一开销，并暴露了紧循环值掩盖的两项真实代价。候选集合由 483 扩到 12,000 时，未激活态的单次边际开销由 0.51 ns 升到 1.19 ns，升幅来自指令缓存与分支目标缓冲：12,000 个一级跳板约占二百 KB 常热代码，L1 指令缓存缺失率由 0.9\% 升到 5.8\%，分支预测失败率由 0.6\% 升到 2.4\%，这两项在 483 个跳板的紧循环里测不到。即便如此，1.19 ns 仍不到激活态 18.79 ns 的十六分之一，\lstinline[breaklines=false]@lat_syscall@ 与 \lstinline[breaklines=false]@lat_ctx@ 的端到端增幅在全内核候选下也只有 2.9\% 与 3.3\%，与第 3 章 483 函数全激活时的开销同量级，而此时随机化的安全覆盖面已扩展到全部内建代码。

\begin{table}[!htbp]
    \bicaption{\enspace 未激活与激活状态在真实规模下的单次调用开销}{\enspace Per-Call Cost in the Inactive and Active States at Realistic Scale}
    \label{tab:4-6}
    \centering
    \scriptsize
    \setlength{\tabcolsep}{3pt}
    \renewcommand{\arraystretch}{1.25}
    \hyphenpenalty=10000\exhyphenpenalty=10000\emergencystretch=1em
    \begin{tabular}{>{\raggedleft\arraybackslash}p{\dimexpr 0.0872\textwidth-2\tabcolsep\relax}>{\raggedleft\arraybackslash}p{\dimexpr 0.1395\textwidth-2\tabcolsep\relax}>{\raggedleft\arraybackslash}p{\dimexpr 0.1220\textwidth-2\tabcolsep\relax}>{\raggedleft\arraybackslash}p{\dimexpr 0.1046\textwidth-2\tabcolsep\relax}>{\raggedleft\arraybackslash}p{\dimexpr 0.1133\textwidth-2\tabcolsep\relax}>{\raggedleft\arraybackslash}p{\dimexpr 0.1395\textwidth-2\tabcolsep\relax}>{\raggedleft\arraybackslash}p{\dimexpr 0.1569\textwidth-2\tabcolsep\relax}>{\raggedleft\arraybackslash}p{\dimexpr 0.1220\textwidth-2\tabcolsep\relax}}
        \hline
        N & 未激活边际开销 & 激活边际开销 & L1I 缺失率 & iTLB 缺失率 & 分支预测失败率 & \lstinline[breaklines=false]@lat_syscall@ 变化 & \lstinline[breaklines=false]@lat_ctx@ 变化\\
        \hline
        483 & 0.51 ns & 18.02 ns & 0.9\% & 0.4\% & 0.6\% & +2.1\% & +2.4\%\\
        3,285 & 0.77 ns & 18.31 ns & 2.3\% & 1.1\% & 1.2\% & +2.4\% & +2.7\%\\
        12,000 & 1.19 ns & 18.79 ns & 5.8\% & 3.7\% & 2.4\% & +2.9\% & +3.3\%\\
        \hline
    \end{tabular}
\end{table}

\section{安全性分析}

\textbf{威胁模型。} 攻击者具备任意内核地址读写能力，目标是获知随机化区域内代码的当前地址并在其失效之前完成控制流劫持。系统已部署持续随机化的执行机制（本文原型见第 3 章），随机化区域的当前地址只存在于跳转桩区的分支立即数中。非随机化区域的控制流完整性由既有内核 CFI 提供（kCFI\citep{clang2022kcfi}、FineIBT\citep{gaidis2023fineibt}、PAL\citep{yoo2022pal}；本文原型见第 5 章）。

\textbf{本章的安全保证。} 其一，当前地址只存在于受硬件保护的位置：跳转桩区处于只执行内存之下，任何数据读取都产生硬件事件，不存在不产生事件的读取路径。其二，读取即触发：事件处理把被读函数交给随机化线程，该函数被激活并搬移，攻击者读到的值在切换完成后失效，失效之后不存在把旧地址翻译回新位置的路径，使用旧地址的控制流落到被覆盖的软件断点并被终止。其三，从事件到失效的窗口有实测的上界：第 4.5.3 节给出该窗口的分布，基础负载下 99 分位不超过 110 μs；第 4.5.4 节的竞态实验进一步给出这一窗口的安全含义，当读取至解引用的时延达到 92 μs 时，攻击成功率已降至 1.0\%。这三条合起来说明，攻击者可以读到当前地址，但可利用它的时间以失效窗口为上界。

\textbf{未激活函数的暴露面。} 未激活函数的当前地址同样只存在于跳转桩中，读取它同样产生事件，但它在被读取之前停在一个静态地址上，读取之后要经过激活与搬移才失效。该函数在被读取之前的地址本就静态，读取之后激活并未制造新的暴露，而是给这段暴露装上了终点。因此未激活函数的失效窗口由两段构成：从读取到跳转桩指向新变体的重定向延迟，以及从激活到旧副本填陷阱的撤销延迟，二者相加才是完整的窗口。

\section{本章小结}

本章把随机化范围由编译期静态选定改为由泄露信号驱动的运行时变量。随机化区域的当前地址集中在跳转桩的分支立即数这一处，对它施加只执行内存，把它当作数据读取的行为就成为可观测事件；正常执行只取指，内核的合法读取者读的是函数体的指令而非跳转桩，机制因此不需要维护任何豁免清单。事件驱动被读函数进入随机化，触发判据与信息泄露之间不再有任何代理。

跳板分为两级之后，未被读取的函数只付一次间接转移的代价，约 5 个周期，被读取的函数才激活为完整的跳板，约 54 个周期。常驻开销由此与候选集合的规模解耦，候选集合可以扩到全内核，而每次调用只为激活集合中的函数付完整的代价。激活是一次 4 字节的原子写入，旧副本的回收由返回标记完成大部分排空，残余以一次性扫描判定。

实验表明，单次检测开销在 ARM 观察点路径上为 1.32 μs，在 x86 保护键路径上为 0.96 μs；基础负载下内核自身的读取每秒触发 0.02 至 0.44 次，开启动态跟踪设施后上升到每秒 0.90 至 6.80 次，其中 ftrace 与 kprobes 分别为 6.80 与 4.20 次；从检测到随机化完成的延迟在基础负载下中位为 12 至 21 μs、99 分位不超过 110 μs。竞态实验表明，当读取至解引用的时延达到 92 μs 时，攻击成功率降至 1.0\%，达到 150 μs 时降至 0.1\%。

本章的机制使随机化区域随泄露事件伸缩，随机化区域之外始终存在布局不变的非随机化区域，其中的代码如何避免被复用，由控制流完整性机制回答（本文原型见第 5 章）。

